% Auto-generated from meeting_2566-12-15_curriculum-structure-review.md (กบ.วช. format v2)
\documentclass{meetingminutes}
\meetingcommittee{คณะกรรมการบริหารหลักสูตร วิศวกรรมคอมพิวเตอร์และปัญญาประดิษฐ์ (บัณฑิตศึกษา)}
\meetingnumber{๔}{๒๕๖๖}
\meetingdate{วันที่ ๑๕ ธันวาคม พ.ศ. ๒๕๖๖}
\meetingplace{ห้องประชุมคณะเทคโนโลยีอุตสาหกรรม}
\meetingstart{๐๙.๐๐ น.}
\meetingend{๑๑.๑๐ น.}

\begin{document}
\maketitle

\psection{ผู้มาประชุม}
\begin{participants}
  \pmember{1}{รองศาสตราจารย์ ดร.อิสระ อินจันทร์}{ที่ปรึกษาหลักสูตร}{ประธาน}
  \pmember{2}{รองศาสตราจารย์ ดร.กันต์ อินทุวงศ์}{ที่ปรึกษาหลักสูตร}{รองประธาน}
  \pmember{3}{ผู้ช่วยศาสตราจารย์ ดร.พิศิษฐ์ นาคใจ}{อาจารย์ประจำหลักสูตร}{กรรมการ}
  \pmember{4}{ผู้ช่วยศาสตราจารย์ ดร.กาญจนา ดาวเด่น}{อาจารย์ประจำหลักสูตร (เลขานุการ)}{กรรมการ/เลขานุการ}
  \pmember{5}{ผู้ช่วยศาสตราจารย์ ดร.พัชรี มณีรัตน์}{อาจารย์ประจำหลักสูตร}{กรรมการ}
  \pmember{6}{ผู้ช่วยศาสตราจารย์ ดร.พิทักษ์ คล้ายชม}{รองคณบดีฝ่ายวิชาการ คณะเทคโนโลยีอุตสาหกรรม}{กรรมการ}
  \pmember{7}{ผู้ช่วยศาสตราจารย์ ดร.โสกณ วิริยะรัตนกุล}{อาจารย์ประจำหลักสูตร}{กรรมการ}
  \pmember{8}{ผู้ช่วยศาสตราจารย์ ดร.ธัชชัย อยู่ยิ่ง}{อาจารย์ประจำหลักสูตร}{กรรมการ}
  \pmember{9}{ผู้ช่วยศาสตราจารย์ ดร.อภิศักดิ์ พรหมฝ่าย}{อาจารย์ประจำหลักสูตร}{กรรมการ}
\end{participants}

\psection{ผู้ไม่มาประชุม}
\noindent ไม่มี\par

\startmeeting
\openingchair{รองศาสตราจารย์ ดร.อิสระ อินจันทร์}{กล่าวเปิดการประชุมและดำเนินการประชุมตามระเบียบวาระ ดังนี้}

\agenda{๑}{รับรองรายงานการประชุมครั้งที่ 3/2566}
ที่ประชุมรับรองรายงานการประชุมครั้งที่ 3/2566 (20 พฤศจิกายน 2566) โดยไม่มีการแก้ไข


\agenda{๒}{รายงาน PLOs ฉบับสมบูรณ์และ Needs-to-PLO Mapping}
ผศ.ดร.พิศิษฐ์ นำเสนอ PLOs ฉบับสมบูรณ์ที่ปรับปรุงจากผลการรับฟังความคิดเห็น พร้อมตาราง Needs-to-PLO Mapping ที่แสดงให้เห็นว่าความต้องการของผู้มีส่วนได้ส่วนเสียสะท้อนอย่างไรในแต่ละ PLO

ที่ประชุมรับทราบและเห็นชอบ PLOs ทั้ง 5 ข้อ พร้อมยืนยันว่าสอดคล้องกับมาตรฐานคุณวุฒิ กมอ. 2565 (4 ด้าน: ความรู้/ทักษะ/จริยธรรม/ลักษณะบุคคล)


\agenda{๓}{ทบทวนโครงสร้างรายวิชาและ Course Description}
ที่ประชุมพิจารณาโครงสร้างรายวิชาทั้งหมด ได้แก่:

วิชาเฉพาะด้านบังคับ (6 วิชา, 18 หน่วยกิต): 1. ปัญญาประดิษฐ์และการเรียนรู้ของเครื่อง (4095101) — 3(2-2-5) 2. การเขียนโปรแกรมขั้นสูงสำหรับการเรียนรู้ของเครื่อง (4095102) — 3(2-2-5) 3. ระเบียบวิธีวิจัยทางวิทยาศาสตร์และวิศวกรรมศาสตร์ (7015906) — 3(2-2-5) 4. สัมมนาวิศวกรรมคอมพิวเตอร์และปัญญาประดิษฐ์ (7015907) — 3(1-4-4) 5. เทคโนโลยีอุบัติใหม่ทางวิศวกรรมคอมพิวเตอร์และปัญญาประดิษฐ์ (7015101) — 3(3-0-6) 6. หัวข้อพิเศษทางวิศวกรรมคอมพิวเตอร์และปัญญาประดิษฐ์ (7015908) — 3(3-0-6)

วิชาเสริม (ไม่นับหน่วยกิต, 6 หน่วยกิต): 1. ภาษาอังกฤษสำหรับนักศึกษาบัณฑิตศึกษา (1555101) 2. การเขียนวิทยานิพนธ์หรือสารนิพนธ์ในระดับบัณฑิตศึกษา (1506001)

ที่ประชุมอภิปรายและเห็นชอบโครงสร้างรายวิชาทั้งหมด


\agenda{๔}{ร่างนโยบายการจัดการเรียนการสอน}
ที่ประชุมเห็นชอบนโยบายการจัดการเรียนการสอนตามปรัชญา IEL (Integrative Experience Learning) ของมหาวิทยาลัย โดย: - เน้นการเรียนรู้เชิงรุก (Active Learning) - บูรณาการประสบการณ์จริงในการเรียนการสอน - เชื่อมโยงทฤษฎีกับการปฏิบัติในบริบทท้องถิ่น

\resolution{เห็นชอบร่างหลักสูตรฉบับสมบูรณ์ในหลักการ และมอบหมายให้คณะทำงานจัดทำ มคอ.2 ฉบับสมบูรณ์เพื่อเสนอต่อคณะกรรมการพิจารณาในปี 2567}

\adjournmeeting

\begin{sigblock}
\typist{ผู้ช่วยศาสตราจารย์ ดร.กาญจนา ดาวเด่น}
\reviewer{รองศาสตราจารย์ ดร.อิสระ อินจันทร์}{กรรมการ}
\chairsig{รองศาสตราจารย์ ดร.อิสระ อินจันทร์}{ประธานการประชุม}
\end{sigblock}

\end{document}
